Let $\groupgen(\secparam)$ be a function that on input $\secparam$ outputs the description of a group $\GG$ with prime order $q$ and generator $g$.
Let the variable $\dl(\groupgen, \adv, \secpar)$ denote the result of the following probabilistic experiment:

\procedureblock{$\dl(\groupgen, \adv, \secpar)$}{
	(\GG, q, g) \gets \groupgen(\secpar) \\
	\text{Sample } s \leftarrow \ZZ_q \\
	s' \leftarrow \adv(\GG, q, g, g^s) \\
	\pcreturn s' = s
}

The $\dl(\groupgen, \adv, \secpar)$ advantage of $\adv$ is defined as
$$
	\advantage{\dl}{\groupgen, \adv} = \prob{\dl(\groupgen, \adv, \secpar) = 1}
$$

If for all adversaries $\adv$, there exists a negligible function $\negl$ such that $\advantage{\dl}{\groupgen, \adv} \leq \negl$, we say that discrete log is hard relative to $\groupgen$.
We may ocasionally abuse notation and say that discrete log is hard for $\GG$ instead, where $\GG$ is representative of the family of groups output by $\groupgen$.
